% Part: proof-theory
% Chapter: cut-elimination
% Section: introduction

\documentclass[../../../include/open-logic-section]{subfiles}

\begin{document}

\olfileid[pt]{pt}{cut}{int}

\olsection{Introdução}

A regra \Cut{} é a regra
\[
\Axiom$\Gamma \fCenter \Delta, !A$
\Axiom$!A, \Pi \fCenter \Lambda$
\RightLabel{\Cut}
\BinaryInf$\Gamma, \Pi \fCenter \Delta, \Lambda$
\DisplayProof
\]
(A !!{formula}~$!A$ é chamada de \emph{fórmula de corte}.)

\Cut{} estava incluída no cálculo de sequentes original de
Gentzen~\Log{LK}. O \emph{Hauptsatz} de Gentzen (``teorema principal'')
afirma que todo sequente !!{provable} em~\Log{LK} também é !!{provable}
sem o uso da regra \Cut{}, O.s.\@, que \Cut{} pode ser ``eliminada'' das
demonstrações de~\Log{LK}. Nossos sistemas \Log{G1c} e \Log{G3c} não
contêm \Cut, mas também se pode levantar a questão para eles: é possível
eliminar \Cut{} de !!{proof}s que usam \Cut{} juntamente com as outras
regras de \Log{G1c} (ou as de \Log{G3c})? Em nossa terminologia já
estabelecida, pode-se perguntar de modo equivalente: \Cut{} é uma regra
admissível de \Log{G1c} (ou \Log{G3c})?

O teorema da eliminação do corte talvez seja o resultado mais importante
e célebre da teoria da prova. Antes de tudo, ele estabelece que
\Log{G1c} e \Log{G3c} não precisam de uma regra que, à primeira vista,
poderia parecer claramente necessária. Pois, ao tentar formalizar
demonstrações efetivas, você perceberá que precisa de uma maneira de
lidar com o uso de lemas. Os matemáticos gostam de demonstrar resultados
e depois recorrer a eles nas demonstrações de outros resultados. Assim,
você poderia primeiro formalizar uma demonstração do resultado $L$ (o
lema) como o !!{proof} de um sequente $\Sequent L$. Em seguida,
formalizaria a demonstração do segundo resultado~$R$, cuja demonstração
usa~$L$, como o !!{proof} de um sequente $L \Sequent R$. Como saber que
existe uma demonstração de~$R$, O.s.\@, !!a{proof} apenas de $\Sequent
R$? É necessário um modo de combinar os dois !!{proof}s, e a regra
\Cut{} permitiria fazer isso.

Mencionamos que \Log{G1c} e \Log{G3c} são completos e, portanto,
demonstram qualquer sequente que se esperaria poder demonstrar---todos os
sequentes válidos. Isso pode ser demonstrado diretamente. Contudo, quando
Gentzen escrevia, sabia-se apenas que os sistemas axiomáticos de derivação
eram completos. Para mostrar que \Log{LK} era completo, Gentzen apresentou
uma tradução de !!{derivation}s do sistema axiomático em
\Log{LK}-!!{proof}s. Essa tradução usava \Cut{} de maneira essencial. A
eliminação do corte não é um resultado importante apenas para a lógica
clássica: muitas outras lógicas têm cálculos de sequentes. Para algumas
lógicas, uma demonstração direta da completude de seus cálculos de
sequentes é difícil ou impossível, de modo que traduções de outros sistemas
(usando \Cut) são a única maneira de mostrar a completude. Nesse caso,
demonstrar a eliminação do corte por métodos da teoria da prova é necessário
para estabelecer que \Cut{} é supérflua e que o sistema sem \Cut{} é
completo.

A eliminação do corte também é importante porque pode ser aplicada para
obter outros resultados de interesse independente. Dois casos que
consideraremos são o lema de interpolação de Craig e o teorema de Herbrand.

As ideias usadas para demonstrar a eliminação do corte normalmente
envolvem mostrar como !!a{proof} que usa \Cut{} pode ser transformada em
uma que não a usa. Essas transformações costumam ser ``locais'', O.s.\@,
tomamos uma parte de !!
a{proof} (geralmente uma sub-!!{proof} que termina em uma inferência
\Cut{}) e a substituímos por outra sub-!!{proof} do mesmo sequente, mas na
qual a inferência \Cut{} foi removida ou substituída por outras. Tais
argumentos fornecem algoritmos passo a passo para transformar !!{proof}s
com \Cut{} em demonstrações sem ela. Isso fornece mais informações quando
procedimentos de eliminação do corte são usados em aplicações. Por
exemplo, há demonstrações do lema de interpolação por métodos da teoria de
modelos, da forma ``se $!A \lif !B$ é válida, então existe um interpolante''
(não importa, por enquanto, o que seja um interpolante). Um argumento da
teoria da prova que usa a eliminação do corte fornece um algoritmo que
recebe !!a{proof} de $!A \lif !B$ e devolve um interpolante. Ao contrário
do argumento da teoria de modelos, o argumento da teoria da prova é
\emph{construtivo}.

Há duas maneiras conhecidas de demonstrar o teorema da eliminação do
corte. Uma delas (que remonta a Gentzen) mostra que é possível remover os
cortes superiores de uma demonstração e então simplesmente os remove até
que nenhum reste. A outra (originalmente devida a Sch\"utte e Tait)
concentra-se nos cortes mais complexos de !!a{proof} e remove
sucessivamente cortes que são máximos no sentido de que nenhum outro corte
tem uma !!{formula} de corte de profundidade maior. Cada uma tem suas
vantagens e desvantagens. Para alguns sistemas, uma abordagem funciona,
enquanto a outra é difícil de implementar ou até mesmo impossível. Por
isso, descrevemos as duas abordagens. Demonstraremos a eliminação do corte
para~\Log{G3c}. Na verdade, mostraremos não que \Cut{} pode ser eliminada,
mas que o \emph{corte com compartilhamento de contexto}~\CutCS pode sê-lo:
\[
\Axiom$\Gamma \fCenter \Delta, !A$
\Axiom$!A, \Gamma \fCenter \Delta$
\RightLabel{\Cut}
\BinaryInf$\Gamma \fCenter \Delta$
\DisplayProof
\]
Se um cálculo de sequentes permite enfraquecimento e contração (seja como
regras propriamente ditas, como em \Log{G1c}, seja como regras admissíveis,
como em \Log{G3c}), \Cut{} pode simular \CutCS, e vice-versa. Escolhemos
\CutCS{} em vez de \Cut porque a primeira estratégia de eliminação do corte
é mais fácil de descrever para~\CutCS, e a segunda estratégia só funciona
para~\CutCS.

\end{document}
